\documentclass{article}

% Packages
\usepackage{amsmath, amsfonts, amssymb, amsthm, bm, bbm} 
\usepackage{enumerate, graphicx, xcolor, tcolorbox, hyphenat, circledsteps}
\usepackage[shortlabels]{enumitem}
\usepackage{url, hyperref}
\hypersetup{
	colorlinks   = true,
	urlcolor     = cyan,
	linkcolor    = red,
	citecolor    = black
}

% Formatting
\usepackage[a4paper, noheadfoot, scale={0.7,0.8}]{geometry}
\title{Wk 4 Supervision Problems 1}
\author{Luca Bollini}
\date{27 October 2026}
\hyphenpenalty=10000

% Commands
\newcommand{\N}{\mathbb{N}}
\newcommand{\Z}{\mathbb{Z}}
\newcommand{\Q}{\mathbb{Q}}
\newcommand{\R}{\mathbb{R}}
\newcommand{\ddt}[1]{\frac{d}{d#1}}
\newcommand{\ddttwo}[1]{\frac{d^2}{d#1^2}}
\newcommand{\ts}{\thinspace}

% Colours
\newtcolorbox{questionbox}{
	colframe=black, 
	colback=white, 
	boxrule=0.2mm,
	arc=0mm,
}

\begin{document}
\maketitle

\begin{center}
\huge{\textbf{Functions}}
\end{center}

\begin{enumerate}
    \item If $A$ is a set with two elements, and $B$ is a set with three elements, how many injective, surjective, and bijective functions are there from $A$ to $B$? How many of each from $B$ to $A$?
    \item Let $f:X\to Y$ be an injective function. Use $f$ to construct an injective function $F:~\mathcal{P}(X) \to \mathcal{P}(Y)$ ($\mathcal{P}$ denotes the power set), justifying why your function is injective.
    \item Let $f:X\to Y$ be a surjective function. Write down a set $X' \subset X$ such that $f|_{X'}$ is bijective. \textit{(You may assume the axiom of choice)}
\end{enumerate}
\begin{questionbox}
    \underline{Axiom of Choice reminder:}\\
    For an arbitrary collection of non-empty sets $\{X_i\}_{i \in I}$, it is possible to select exactly one element from each set as a representative. That is, there exists $\{x_i\}_{i\in I}$ such that $x_i \in X_i$.
\end{questionbox}

\begin{center}
\huge{\textbf{Fixed Points of Difference Equations}}
\end{center}

\begin{enumerate}
    \item Consider the difference equation
    \[
    x_{n+1} = \lambda x_n - x_n^3,
    \]
    where $\lambda$ is a non-negative parameter.
    Find the fixed points of this equation and determine their stability for different values of $\lambda$.
    \item Let $f: \R\to\R$ be a continuous function, and consider the difference equation
    \[
    x_{n+1} = f(x_n)
    \]
    Prove (informally) that if this equation has a 2-cycle with the values $a$ and $b$, then there exists a fixed point $x^* \in (a,b)$.
\end{enumerate}

\end{document}